\chapter{GPU architecture: A100 and H100}\label{ch:gpuarch}

A GPU is a processor built for throughput: it runs tens of thousands of threads at once and hides the latency of
each memory access behind the work of other threads. This chapter explains the parts of an NVIDIA GPU that matter
for WRF, and compares the two GPUs the port targets:
\begin{itemize}
  \item the A100 (Ampere, compute capability 8.0) \citep{nvidia2020a100};
  \item the H100 (Hopper, compute capability 9.0) \citep{nvidia2022h100}.
\end{itemize}
The numbers in this chapter are datasheet values. The port measures the real ceilings of its machines with
micro-benchmarks (\cref{ch:prof}), and the performance models of \cref{ch:perf} use the measured values.

\section{Throughput instead of latency}

A CPU core is built to finish one instruction stream as fast as possible. It has large caches, branch prediction
and out-of-order execution, all to keep the latency of each instruction low. A GPU spends the same transistors on
many simple arithmetic units and on registers for many threads. A single GPU thread is slow; a GPU is fast only when
it has enough independent threads to keep its units busy.

WRF's dynamics fits this model well: almost every loop nest updates each grid point of a field independently of the
others. The d02 domain of the case has $810 \times 810 \times 59 \approx 3.9 \times 10^7$ mass points, so a
pointwise loop offers almost forty million independent iterations. Column physics offers fewer, one per column
($6.6\times10^5$ on d02), each with a long sequential body. The fire grid offers $10^7$ nodes.

\section{Streaming multiprocessors, warps and threads}

An NVIDIA GPU consists of \emph{streaming multiprocessors} (SMs): 108 on the A100, 132 on the H100 SXM and 114 on
the H100 PCIe. A program launches a \emph{kernel} as a grid of \emph{thread blocks}. Each block runs on one SM, and
an SM holds several blocks at once.

\paragraph{Warps.} The threads of a block are executed in groups of 32, called \emph{warps}. The 32 threads of a
warp execute the same instruction at the same time on different data (single instruction, multiple threads, SIMT).
Each SM has four sub-partitions, each with its own warp scheduler, which issues one instruction per clock from one
of its ready warps.

\paragraph{Latency hiding.} When a warp waits for memory, which takes hundreds of clock cycles, the scheduler simply
issues from another warp. An SM can hold up to 64 warps (2048 threads). The more warps are resident, the better the
latency is hidden. The fraction of the 64 slots in use is the \emph{occupancy}.

How many accesses must be in flight follows from Little's law: to sustain a bandwidth $B$ with a memory latency
$L$, the GPU must have $B \cdot L$ bytes in flight. With $B = 1.56$~TB/s and an assumed latency of 500~ns, that is
about 780~KB, or 7~KB per SM of the A100: some 1800 outstanding 4-byte loads per SM. Only many resident warps, each
with several independent loads, supply that.

\paragraph{Divergence.} When the threads of a warp take different branches of an \code{IF}, the warp executes both
branches one after the other, with the threads of the other branch masked off. The cost is the sum of both branches.
WRF has many such branches: the fire's WENO/ENO choice, the boundary rows of advection, the night columns of the
shortwave scheme (\cref{ch:warp}).

\section{The memory hierarchy}

\Cref{fig:memhier} shows the levels of memory a kernel sees, from the fastest and smallest to the slowest and
largest.

\begin{figure}[ht]
\centering
\begin{tikzpicture}[lvl/.style={draw, rounded corners, minimum height=0.9cm, align=center, font=\small}]
  \node[lvl, fill=green!10, minimum width=4cm] (r) at (0,4.4) {Registers: 64\,K $\times$ 32 bit per SM\\(at most 255 per thread)};
  \node[lvl, fill=green!6, minimum width=6cm] (s) at (0,3.2) {L1 cache and shared memory, per SM\\A100: 192 KB ($\le$ 164 KB shared); H100: 256 KB ($\le$ 228 KB shared)};
  \node[lvl, fill=yellow!10, minimum width=8cm] (l2) at (0,2.0) {L2 cache, shared by all SMs\\A100: 40 MB; H100: 50 MB};
  \node[lvl, fill=orange!10, minimum width=10cm] (m) at (0,0.8) {Device memory (HBM): 40 or 80 GB\\A100 40 GB: 1.56 TB/s; A100 80 GB SXM: 2.04 TB/s; H100 PCIe: 2.0 TB/s; H100 SXM: 3.35 TB/s};
  \node[lvl, fill=red!6, minimum width=12cm] (h) at (0,-0.4) {Host memory, over PCIe or NVLink: tens of GB/s};
  \draw[-{Latex}] (6.6,4.6) -- node[right, font=\small, align=left] {larger,\\slower} (6.6,-0.5);
\end{tikzpicture}
\caption{The memory hierarchy of the A100 and H100 (datasheet values). Measured bandwidths and latencies of the
port's machines are \tbr\ (profiler tasks PR-H1 to PR-H4).}
\label{fig:memhier}
\end{figure}

\paragraph{Registers.} Each thread's scalars live in registers. The SM's register file of 65\,536 32-bit registers
is shared by all resident threads, so the registers per thread limit the occupancy:
\begin{equation}
  \text{resident warps per SM} \le \min\left(64,\ \left\lfloor \frac{65\,536}{32\,R} \right\rfloor\right),
  \label{eq:occ-regs}
\end{equation}
where $R$ is the number of registers per thread, rounded up to the allocation granularity. A kernel with $R = 32$
can fill all 64 warps; one with $R = 128$ only 16, an occupancy of 25\,\%.

\paragraph{Local memory.} What does not fit in registers is \emph{spilled} to local memory: a per-thread area in
device memory, cached in L1 and L2. Private arrays indexed by a loop variable also go there, because registers
cannot be indexed. Column physics, with dozens of private column arrays per thread, lives largely in local memory
(\cref{ch:colphys}).

\paragraph{Shared memory and L1.} Each SM has a fast on-chip memory that is split between the L1 cache and
\emph{shared memory}, which a thread block manages explicitly. Shared memory lets the threads of a block load a tile
of an array once and read it many times: the basis of tiling (\cref{ch:tiling}).

\paragraph{L2 cache.} All SMs share the L2 cache. Data that a kernel writes and the next kernel reads may still be in
L2 if it fits:
\begin{itemize}
  \item a dev-case fire field ($724^2$ values, 2~MB) fits easily;
  \item a d02 3D field ($821 \times 60 \times 821$ values, 162~MB) does not;
  \item a full-case fire field ($3284^2$ values, 43~MB) lies between the two GPUs: just over the A100's 40~MB and
    just under the H100's 50~MB.
\end{itemize}

\paragraph{Device memory.} High-bandwidth memory (HBM2e on the A100, HBM3 on the H100 SXM) holds all of the model's
state. Its bandwidth sets the speed of most of WRF. Reading one d02 3D field once takes
\begin{equation}
  \frac{162\ \text{MB}}{1.56\ \text{TB/s}} \approx 104\ \mu\text{s on an A100 40 GB}, \qquad
  \frac{162\ \text{MB}}{3.35\ \text{TB/s}} \approx 48\ \mu\text{s on an H100 SXM}.
\end{equation}

\paragraph{Host memory.} The host is reached over PCIe (A100 PCIe: Gen4, about 32~GB/s per direction; H100 PCIe:
Gen5, about 64~GB/s) or NVLink. This is two orders of magnitude slower than device memory. The port therefore keeps
the whole model state on the device and copies only what the host must see (\cref{ch:port}).

\section{Coalescing}

Device memory is accessed in 32-byte \emph{sectors}. When the 32 threads of a warp load 32 consecutive 4-byte
values, the warp needs exactly four sectors: the access is \emph{coalesced}. When the threads load values that are
far apart, each needs its own sector, and the warp moves eight times more data than it uses.

WRF stores 3D fields as \code{a(i,k,j)}, with $i$ the fastest index in memory. A kernel whose consecutive threads
take consecutive values of $i$ is coalesced. The port's kernels keep the source's loop order (\code{j}, \code{k},
\code{i}, with \code{i} innermost), and the compiler gives consecutive $i$ to consecutive threads. A column kernel,
with one thread per $(i,j)$ and a sequential loop over $k$ inside, is coalesced too: at each $k$ the threads of a warp
read consecutive $i$ (\cref{ch:memory}).

\section{Arithmetic}

Each SM has units for single precision (FP32), double precision (FP64) and integers.
\begin{itemize}
  \item The A100 peaks at 19.5~TFLOP/s in FP32 and 9.7~TFLOP/s in FP64; the H100 SXM at about 67 and 34.
  \item These peaks count a fused multiply-add (FMA) as two operations. The port's REPRO mode forbids FMA
    contraction (\cref{ch:repro}), so its peak is half of these numbers.
  \item WRF computes in FP32. Its double-precision work is the inside of the reproducible elementary functions
    (\code{rp_*}), which evaluate in FP64 and round once to FP32.
\end{itemize}
For most of WRF the arithmetic peak does not matter. A stencil kernel does a few operations per value it loads, far
below the ratio of compute to bandwidth of either GPU (\cref{ch:perf}). Column physics and the reproducible
transcendental functions are the exceptions.

\section{A100 and H100 compared}

\begin{table}[ht]
\centering
\small
\begin{tabular}{lcccc}
\toprule
 & A100 40 GB & A100 80 GB SXM & H100 PCIe 80 GB & H100 SXM 80 GB \\
\midrule
Compute capability & 8.0 & 8.0 & 9.0 & 9.0 \\
SMs & 108 & 108 & 114 & 132 \\
HBM bandwidth (TB/s) & 1.56 & 2.04 & 2.0 & 3.35 \\
L2 cache (MB) & 40 & 40 & 50 & 50 \\
L1 + shared per SM (KB) & 192 & 192 & 256 & 256 \\
Max shared per block (KB) & 164 & 164 & 228 & 228 \\
FP32 / FP64 (TFLOP/s) & 19.5 / 9.7 & 19.5 / 9.7 & $\approx$ 51 / 26 & $\approx$ 67 / 34 \\
Registers per SM & 64 K & 64 K & 64 K & 64 K \\
Max threads per SM & 2048 & 2048 & 2048 & 2048 \\
\bottomrule
\end{tabular}
\caption{The GPUs of the port (datasheet values; \citealp{nvidia2020a100,nvidia2022h100}).}
\label{tab:gpus}
\end{table}

Beyond the numbers in \cref{tab:gpus}, the H100 adds features that change how a kernel can be written:
\begin{description}
  \item[Tensor Memory Accelerator (TMA).] A unit that copies a multi-dimensional tile between device memory and
    shared memory asynchronously, given only the tile's coordinates. A stencil kernel can load its next tile while it
    computes on the current one, without spending threads on address arithmetic.
  \item[Thread-block clusters.] Up to 16 blocks on neighbouring SMs form a cluster. They can read each other's shared
    memory (\emph{distributed shared memory}), so neighbouring tiles of a stencil can share their halos.
  \item[Larger L2 and shared memory.] 50 instead of 40~MB of L2, and 228 instead of 164~KB of shared memory per
    block. This changes the best tile size (\cref{ch:tiling}).
\end{description}
The A100 already has asynchronous copies from device to shared memory, and a control of the L2 cache that keeps a
chosen address range resident (\emph{persisting} accesses, \cref{ch:memory}).

\section{What to expect for WRF}

From these facts, and before any measurement:
\begin{itemize}
  \item \textbf{Memory-bound stencils} (most of the dynamics) should run in proportion to the memory bandwidth:
    about twice as fast on an H100 SXM as on an A100 40 GB.
  \item \textbf{Short kernels} gain less. The acoustic loop launches many kernels that each run for a few
    microseconds, and the cost of a launch does not shrink with bandwidth. Fewer launches and fused kernels matter
    more on the H100 (\cref{ch:fusion}).
  \item \textbf{Column physics} is limited by registers, local memory and divergence rather than bandwidth. The
    larger register and L1 budget of the H100 changes the best batch shape (\cref{ch:colphys}).
  \item \textbf{The fire grid} of the full case fits the H100's L2 but not the A100's, so the fire kernels may scale
    differently from the rest. Measurements on the dev case, whose fire grid fits both caches, must not be
    extrapolated to the full case.
\end{itemize}

\begin{portnote}
The port builds one executable for both GPUs (\texttt{-gpu=cc80,cc90}). Both must give the same bits as CPU-REF,
and therefore as each other (G5 item 5). The performance settings that differ between the GPUs, such as vector
lengths and register limits, change only the schedule, never the arithmetic. They are kept per GPU in
\code{port/ENVIRONMENT.md}.
\end{portnote}
